\documentclass[10pt,a4paper]{amsart}
\usepackage[margin=19mm]{geometry}
\usepackage{amsmath,amssymb,mathtools}
\usepackage[scr=rsfs,cal=euler]{mathalfa}
\usepackage{xfrac}
\usepackage[unicode,hidelinks,bookmarks=false]{hyperref}
\newcommand{\ie}{i.e.\ }
\newcommand{\cf}{cf.\ }
\newcommand{\resp}{respectively\ }
\newcommand{\Subsection}{\subsection}
\providecommand{\nobreakdash}{\nobreak\hbox{-}}
\setlength{\emergencystretch}{2em}
\multlinegap0pt
\usepackage[shortlabels]{enumitem}
\usepackage{amssymb}
\usepackage{mathtools}
\newtheorem{theorem}{Theorem}
\newtheorem{proposition}{Proposition}
\newtheorem{lemma}{Lemma}
\newtheorem{corollary}{Corollary}
\usepackage{cleveref}
\crefname{theorem}{Theorem}{Theorems}
\crefname{proposition}{Proposition}{Propositions}
\crefname{lemma}{Lemma}{Lemmas}
\crefname{corollary}{Corollary}{Corollarys}
\newcommand{\R}{\mathbb{R}}
\newcommand{\Sn}{\mathbb{S}}
\newcommand{\Span}{\operatorname{span}}
\newcommand{\cS}{\mathcal{S}}
\title{Rank-Three Projections and Minimal Multiplicity Bipartitions of Path Complements}
\author{Jintao Fei, Jiangying Luo}
\date{Source: arXiv:2608.27227v1 (CC BY 4.0)}
\begin{document}
\maketitle

\noindent\emph{Source and licence.} Rank-Three Projections and Minimal Multiplicity Bipartitions of Path Complements. Version arXiv:2608.27227v1 (CC BY 4.0), distributed by arXiv under the Creative Commons Attribution 4.0 International licence, http://creativecommons.org/licenses/by/4.0/, as recorded in the arXiv licence metadata for this exact version and shown on the abstract page https://arxiv.org/abs/2608.27227v1. Full licence notice: \texttt{licenses/arxiv-2608.27227-CC-BY-4.0.txt}.\par\smallskip

\noindent\textit{Editorial scope.} Counted unit: the article's theorem thm:main (source0.tex lines 127--135). The article's complete original proof of it is delivered with it (source0.tex lines 391--445, one \texttt{proof} environment of 55 lines, which the article places away from the statement). No mathematics is written, repaired or re-derived: the statement and the proof are the article's own lines, in the article's own notation, and no retained line is substituted or re-flowed.  Retained uncounted as the source-local context the proof needs: lines 168--179; lines 200--219; lines 312--316; lines 322--344.  Nothing was omitted from any retained span, so the package carries no omission record.  No retained line contains a citation, so the wrapper carries no bibliography and no bibliography entry.  No retained line contains a figure environment, an image-inclusion command or drawing code, so no image is delivered and the package declares no asset dependency.  Editorial formatting only, in addition: an AMS \texttt{amsart} preamble that restates the article's own declarations and macros used by the retained lines, namely the environments \texttt{theorem}, \texttt{proposition}, \texttt{lemma}, \texttt{corollary} on separate counters, together with the standard packages those lines need.  The retained environments print in this document as Theorem 1, Proposition 1, Lemma 1, Corollary 1 in order of appearance, whereas the article prints the same items with the numbering of its own full document.  The complete native source of the article as distributed in the arXiv e-print for this exact version is bundled unchanged as \texttt{sources/208697/source0.tex}.
\begin{proposition}[Projection--frame equivalence]\label{prop:projection}
Let $G$ be a graph on $n$ vertices and let
$1\le r\le\lfloor n/2\rfloor$.  The bipartition $[n-r,r]$ is achievable by
$G$ if and only if there is an $n\times r$ real matrix $U$ such that
\[
 U^{\mathsf T}U=I_r
 \quad\text{and}\quad
 UU^{\mathsf T}\in\cS(G).
\]
In this case $UU^{\mathsf T}$ is a rank-$r$ orthogonal projection with
spectrum $\{0^{(n-r)},1^{(r)}\}$.
\end{proposition}

\begin{lemma}[Anchor absorption]\label{lem:absorption}
Let $x_1,\ldots,x_r\in\R^d$ be nonzero vectors such that
\[
 \Span\{x_i x_i^{\mathsf T}:1\le i\le r\}=\Sn^d
 \tag{2.1}\label{eq:span}
\]
and suppose that, for some $a_1,\ldots,a_r>0$,
\[
       \sum_{i=1}^r a_i x_i x_i^{\mathsf T}=I_d.
 \tag{2.2}\label{eq:anchorparseval}
\]
Then for every finite collection $y_1,\ldots,y_s\in\R^d$ there are strictly
positive weights $\alpha_1,\ldots,\alpha_r$ and
$\beta_1,\ldots,\beta_s$ satisfying
\[
 \sum_{i=1}^r \alpha_i x_i x_i^{\mathsf T}
 +\sum_{j=1}^s \beta_j y_j y_j^{\mathsf T}=I_d.
 \tag{2.3}\label{eq:absorbed}
\]
\end{lemma}

\begin{corollary}\label{cor:extend}
Every integral faithful path chain $d_1,\ldots,d_k$ in $\R^3$ whose vectors
are pairwise nonparallel extends, for every $N\ge k$, to an integral
faithful path chain $d_1,\ldots,d_N$ with the same property.
\end{corollary}

\begin{proposition}[A six-vector absorbing anchor]\label{prop:anchor}
Let
\[
\begin{aligned}
 d_1&=(2,-1,0),      & d_2&=(0,0,1),       & d_3&=(1,1,0),\\
 d_4&=(1,-1,-2),     & d_5&=(0,2,-1),      & d_6&=(2,1,2),
\end{aligned}
\tag{4.1}\label{eq:seedvectors}
\]
and
\[
 (a_1,\ldots,a_6)
 =\left(\frac18,\frac38,\frac14,\frac1{12},\frac18,\frac1{24}\right).
\tag{4.2}\label{eq:seedweights}
\]
Then:
\begin{enumerate}[label=\textup{(\roman*)}]
\item $d_1,\ldots,d_6$ form a faithful path chain and are pairwise
      nonparallel;
\item $\sum_{i=1}^6 a_i d_i d_i^{\mathsf T}=I_3$;
\item the six matrices $d_i d_i^{\mathsf T}$ form a basis of $\Sn^3$.
\end{enumerate}
\end{proposition}

\begin{theorem}\label{thm:main}
For every integer $n\ge6$,
\[
             MB(\overline{P_n})=3.
\]
Equivalently, $\overline{P_n}$ has a realization that is a rank-three
orthogonal projection, and it has no realization by a projection of rank
one or two.
\end{theorem}

\begin{proof}[Proof of \Cref{thm:main}]
Fix $n\ge6$.  By \Cref{cor:extend}, the six vectors in
\eqref{eq:seedvectors} extend to integral vectors
$d_1,\ldots,d_n$ satisfying
\[
 d_i^{\mathsf T}d_j=0
 \quad\Longleftrightarrow\quad |i-j|=1
 \qquad(i\ne j).
\tag{4.6}\label{eq:fullchain}
\]
If $n=6$ no extension is needed.  Apply \Cref{lem:absorption} with the first
six vectors as anchors and $d_7,\ldots,d_n$ as the finite tail.  By
\Cref{prop:anchor}, there exist weights $w_1,\ldots,w_n>0$ such that
\[
                 \sum_{i=1}^n w_i d_i d_i^{\mathsf T}=I_3.
\tag{4.7}\label{eq:fullparseval}
\]
Let $U$ be the $n\times3$ matrix whose $i$th row is
$u_i^{\mathsf T}=\sqrt{w_i}\,d_i^{\mathsf T}$.  Then
$U^{\mathsf T}U=I_3$, and hence
\[
                  P=UU^{\mathsf T}
\]
is a rank-three orthogonal projection.  For $i\ne j$,
\[
 P_{ij}=\sqrt{w_iw_j}\,d_i^{\mathsf T}d_j.
\]
All weights are positive, so \eqref{eq:fullchain} says precisely that
$P_{ij}=0$ if and only if $|i-j|=1$.  The consecutive pairs are exactly the
nonedges of $\overline{P_n}$, and therefore
$P\in\cS(\overline{P_n})$.  By \Cref{prop:projection},
\[
                  MB(\overline{P_n})\le3.
\]
This construction also proves directly that $q(\overline{P_n})=2$: the
projection has two eigenvalues, while $\overline{P_n}$ has an edge and hence
cannot be represented by a scalar matrix.

For the reverse inequality, suppose that
$P=UU^{\mathsf T}\in\cS(\overline{P_n})$ is a projection of rank
$r\in\{1,2\}$, and let $u_i^{\mathsf T}$ be row $i$ of $U$.  Every $u_i$ is
nonzero, since a zero row would make vertex $i$ isolated in the graph of
$P$, whereas $\overline{P_n}$ has no isolated vertex for $n\ge4$.  Moreover,
\[
                 u_i\perp u_{i+1}\qquad(1\le i<n).
\tag{4.8}\label{eq:roworth}
\]
Rank one is impossible because two nonzero vectors in $\R$ cannot be
orthogonal.  If $r=2$, then $u_1\perp u_2$ and $u_3\perp u_2$ force
$u_1\parallel u_3$.  Since $u_3\perp u_4$, it follows that
$u_1\perp u_4$.  But $\{1,4\}$ is an edge of $\overline{P_n}$, so
$P_{14}$ must be nonzero, a contradiction.  Thus no rank-one or rank-two
projection has the required zero pattern, and
$MB(\overline{P_n})\ge3$.
\end{proof}

\end{document}
